\section{Positive critical marked sources with collapsing clearance slices}
\label{pc:section}

A physical corner can have a positive incoming critical fan although
its angular section at zero clearance has collapsed. Such a fan is
not an angular-null error. We now pay its actual leading coefficient
by transporting finitely many positive physical sectors to a common
nearby positive clearance, and using the branch-supremum theorem there.
This argument concerns the original unvisited level $u=0$ and permits
simultaneous unvisited tangencies and original section cuts. All
visited incidences at a limiting germ remain positive. A trace at
an actual visited grazing collision is a different measure.

Fix $R\in[9/20,47/100]$, $m\ge1$ and $0\le j<m$. A marked
critical germ has the original $m$-flight strict visited word,
normal endpoints $p_0=p_m=0$, and an unvisited marked disk with
a perpendicular interior foot on its $j$th flight at clearance
zero. Its local physical side consists of actual regular records
with $u>0$, all the original other-obstacle inequalities, and the
specified original section and finite record tests. Original
physical roof intervals may also be imposed. Germs with the same
limiting marked state and strict visited word have their physical
side sets united and disjointized, not counted a second time.
All actual marked disks remain separate marked copies.

Use the native normalized coordinates
$y=G_x^{1/2}(q-q_x)$ and the physical angular fraction
$\beta_y(x)\in[0,1]$ from Section~\ref{sec:mc-critical}.
The incoming leading coefficient is
\begin{equation}\label{pc:coefficient}
 \mathfrak J(x)\beta_y(x),\qquad
 \mathfrak J(x)=\frac{\sqrt{\det G_x}}{2R^2c_*}\,u_x,
 \qquad
 u_x=\frac{R|G_{0m}(x)|}{\det G_x}.
\end{equation}
These are the actual auxiliary-action boundary values of
\eqref{eq:mc-weight}. No derivative of the singular physical
collision map is assigned at an unvisited tangency. The physical
map equals this auxiliary map on the actual regular side.

\subsection{A physical representative at every small positive clearance}

\begin{lemma}[Analytic sectors and common one-sided levels]
\label{pc:representatives}
At a fixed marked critical germ as above, its finite physical-side
conditions can be written locally as a finite Boolean combination
of real analytic inequalities, including $u>0$.
The limiting angular indicator exists outside finitely many rays.
If $\beta_y(x)>0$, there are actual regular points $x(u)$ on this
physical side for every sufficiently small $u>0$, with their
designated marked clearance exactly $u$ and $x(u)\to x$.
Representatives may additionally avoid the cuts and exceptional
angular sets of both partitions in Theorem~\ref{thm:mb-branch-sup}. For any finite
collection of such germs they may be chosen at every common
sufficiently small $u>0$.
\end{lemma}

\begin{proof}
Positive visited incidences give an invertible internal optical
Hessian. The local strict-word reflection equations and all contact
positions are therefore real analytic functions of the endpoint
coordinates by the implicit-function theorem. The flight lengths
are bounded away from zero. Fixing the finite word and a sufficiently
small neighborhood, only finitely many lattice centers can meet
one of its flights or its clearance neighborhood. Other centers
have a fixed positive distance and impose no active local cut.
The first-hit and unvisited-obstacle requirements are finite
analytic sign conditions: the perpendicular-foot tests and squared
line-distance tests may be used on their respective finite cases.
Endpoint distance tests are analytic as well. The designated
marked foot is strictly interior and separated from the visited
contacts by the original standoff estimate.

Each of the original finitely many section rectangles is given by
coordinate inequalities pulled back through finite regular maps.
Occupation is a finite Boolean combination of those tests; the
lifted displacement is fixed on the strict word. The physical roof
is the same analytic length sum, so a fixed roof interval adds two
analytic inequalities, including at a critical roof value. No
replacement section or time is used. This reasoning concerns
these original finite analytic tests; it makes no ray-persistence
assertion for an arbitrary additional Borel selector.

In any analytic local coordinate $z$ centered at the germ, remove
constant truth values and identically zero functions from the finite
Boolean formula. For each remaining analytic function $f$ with
$f(0)=0$, let $P_f$ be its first nonzero homogeneous Taylor term.
The nonzero polynomial $P_f$ vanishes on only finitely many rays
of the unit circle. On each complementary open angular interval,
the signs of all $f(re)$ are fixed for sufficiently small $r>0$,
uniformly on a compact smaller interval. The formula thus has a
limiting indicator there. Boundary rays have zero angular measure.
An invertible linear coordinate change preserves positivity of the
angular measure, so $\beta_y>0$ supplies at least one such open
interval of actual physical directions. In the inherited local radial
coarea formula the optical coefficient is continuous and bounded
at this fixed positive-visited germ. Dominated convergence of the
angular indicators therefore gives exactly
$\mathfrak J(x)\beta_y(x)$ as its leading one-sided coefficient.
In particular zero opening contributes zero to this leading
coefficient; this is not an assertion about bulk-null traces or
about a collision-uniform coarea remainder.

Take $z=(p_0,p_m)$. Its local change from $q$ is invertible by the
positive Hessian and nonzero optical corner in
Lemma~\ref{lem:mc-hessian}. The exact half-ray determinant of
Theorem~\ref{thm:bf-foot-determinant}, continued on the same strict
word to the limit, makes $(\vartheta,u)$ another local analytic
coordinate system:
\[
 \left|\det D_{(\vartheta,u)}(p_0,p_m)\right|
             =A_-^pA_+^p(h_-+h_+)>0 .
\]
Consequently $D_zu(0)\ne0$. Since the physical side includes $u>0$,
a positive physical angular interval contains a direction $e$ with
$D_zu(0)e>0$, after omitting the two zero directions of this linear
form. The ray $z=re$ satisfies all selected physical sign conditions
strictly for small $r>0$, after constant truth values are accounted
for. Along it
\[
 u(re)=rD_zu(0)e+O(r^2),\qquad
 \frac{d}{dr}u(re)=D_zu(0)e+O(r)>0 .
\]
The one-variable inverse-function theorem supplies a representative
at every small positive $u$, converging to the germ.

At each such representative the same sign conditions persist on
an open neighborhood. Intersecting with its fixed-$u$ coordinate
line leaves a nonempty open angular interval. The two one-sided
partitions in 319 omit only angular-null sets at that fixed level.
Choose a point in this interval outside their union, at angular
distance tending to zero from the original representative as
$u\downarrow0$. It still satisfies the same physical tests and
still converges to $x$. For a finite family, take the minimum of
the finitely many positive level ranges and perform these choices
for every common $u$ in that range. No common neighborhood size
over all words, radii or collision counts is asserted.
\end{proof}

\subsection{The critical branch-supremum transfer}

\begin{lemma}[Both half-ray sums on positive physical critical germs]
\label{pc:half-sums}
Let $\mathcal F_{m,j}(R)$ be the complete one-copy family of the
specified zero-clearance marked critical germs with
$\beta_y(x)>0$. At their limiting strict-word optical values,
\begin{equation}\label{pc:half-sums-bound}
 \sum_{x\in\mathcal F_{m,j}(R)}(A_-^p(x))^{-1}\le H_{\rm nm},
 \qquad
 \sum_{x\in\mathcal F_{m,j}(R)}(A_+^p(x))^{-1}\le H_{\rm nm}.
\end{equation}
There is no assumption that a germ has a nonempty regular angular
interval at $u=0$.
\end{lemma}

\begin{proof}
First fix any finite set $\mathcal F$ of distinct merged germs.
Translate each designated marked disk to the origin and retain its
impact sign. A limiting angular direction and that sign determine
one marked tangent line. Following its two halves, the next strictly
met disk is unique, even when the limiting line is tangent to other
earlier disks. Continuing through its positive-incidence reflections
determines the complete strict visited word. Using an earlier
tangency as a visited collision would introduce a zero incidence
and would belong to a different, here unasserted stratum.
This is the strict-hit uniqueness argument of
Theorem~\ref{thm:mc-critical-sectors}, applied also backwards to
recover the original initial disk at index $j+1$.

Thus distinct merged germs with the same impact sign have distinct
limiting angular points. If they give the same marked point, they
are the same strict-word germ and their side selectors have already
been united. The original initial-disk translation is unique and
does not add a lattice multiplicity. Choose an angular seam avoiding
the finitely many limiting points. The constants in 319 are uniform
in this initial seam choice.

For every common sufficiently small $u>0$, take representatives
from Lemma~\ref{pc:representatives} inside intervals of both partitions, away from their cuts and
exceptional angular sets. For small enough $u$, no two representatives with the
same sign can lie in the same backward interval
$I\in\mathcal P^-_{j+1}$. Otherwise monotonicity and the pointwise
lower bound in 319 give
\[
 |p_0(x(u))-p_0(\widehat x(u))|
       \ge k_-a_\circ
                  |\vartheta(x(u))-\vartheta(\widehat x(u))|.
\]
The left side tends to zero, whereas the right side stays bounded
below by the separation of two distinct limiting angles.
A finite family has a common positive minimum separation for
each sign, so this contradiction excludes every such pair for
all sufficiently small $u$. The same argument uses $p_m$ and
$\mathcal P^+_{m-j}$ for the forward halves. Both momenta tend to
zero because the limiting endpoints are normal. No derivative of
the selectors is involved.

The representatives therefore occupy distinct backward intervals,
and also distinct forward intervals. Each reciprocal half derivative
is bounded by the supremum of its actual interval. Applying 319
at this common positive $u$ gives
\[
 \sum_{x\in\mathcal F}(A_\pm^p(x(u)))^{-1}\le H_{\rm nm}.
\]
Every visited incidence at each fixed limiting germ is positive,
so its finite strict-word optical matrix depends continuously on
these representatives. Passing $u\downarrow0$ in the finite sum
proves the two inequalities for $\mathcal F$. Taking the supremum
over finite subsets proves \eqref{pc:half-sums-bound}. Equivalently, this
defines the full nonnegative sum without choosing an order; finiteness
implies at most countably many positive terms.

The limit was taken at fixed $R,m$ and finite $\mathcal F$ before
that exhaustion. Uniformity in $R,m$ comes from the same final
constant, not from an interchange of local and long-time limits.
In particular a collapsed angular slice at level zero never had
to be populated by actual regular points.
\end{proof}

\begin{theorem}[Complete positive marked critical-fan payment]
\label{pc:fan}
Put $\sigma=47/59$ and retain $Q_{\rm mb}$ from
\eqref{eq:mb-constants}. For the original zero-clearance marked
physical critical family with positive visited incidences,
\begin{align}
 \sum_{x:\,\beta_y(x)>0}u_x
       &\le Q_{\rm mb}\,
                       \sigma^{\max(j,m-1-j)},                 \label{pc:weights}\\
 \mathfrak F_{m,j}(R):=
       \sum_x\mathfrak J(x)\beta_y(x)
       &\le C_{\rm pc}\,
                       \sigma^{\max(j,m-1-j)},\qquad
 C_{\rm pc}=\frac{170}{18c_*}\,Q_{\rm mb}.                       \label{pc:fan-bound}
\end{align}
Both impact signs, all actual marked disks and strict words, and
all simultaneous finite original physical cuts are included.
The bounds are uniform in radius and collision count and may be
restricted to any original exact record or to critical values in
any interval of the unchanged physical roof. Any further weights
between zero and one can decrease this already bounded positive sum.

In particular the complete sum over \emph{all} marked flight
indices, including shallow indices, satisfies
\begin{equation}\label{pc:all-marks}
 \sum_{j=0}^{m-1}\mathfrak F_{m,j}(R)
        \le C_{\rm pc}m\,\sigma^{(m-1)/2},\qquad
 \sup_R m^2\sum_{j=0}^{m-1}\mathfrak F_{m,j}(R)
        \le C_{\rm pc}m^3\sigma^{(m-1)/2}\longrightarrow0.
\end{equation}
The same conclusion holds after the original exact-label suprema.
No outer width or guard is changed.
\end{theorem}

\begin{proof}
The two actual half counts are $j+1$ and $m-j$. At the limiting
strict-word germ, continuity and the positive optical recurrence
of 319 give
\[
 (A_-^p)^{-1}\le E_*\sigma^j,\qquad
 (A_+^p)^{-1}\le E_*\sigma^{m-1-j}.
\]
The same exact optical determinant gives
\[
 u_x\le C_{\rm foot}(A_-^p)^{-1}(A_+^p)^{-1}.
\]
Bound either factor individually and sum the other by
Lemma~\ref{pc:half-sums}. The two resulting bounds are
$Q_{\rm mb}\sigma^j$ and $Q_{\rm mb}\sigma^{m-1-j}$.
Their minimum is \eqref{pc:weights}, since $0<\sigma<1$.

Use the actual coefficient identity \eqref{pc:coefficient} and
the sharp inherited critical bound
$\sqrt{\det G_x}/R^2\le170/9$ from Lemma~\ref{lem:mc-hessian}.
The united one-copy fraction satisfies $0\le\beta_y\le1$.
Zero-fraction germs contribute zero by their defining angular
coefficient, not by a bulk-null claim. Thus
\[
 \sum_x\mathfrak J(x)\beta_y(x)
 \le\frac{170}{18c_*}
       \sum_{\beta_y(x)>0}u_x,
\]
proving \eqref{pc:fan-bound}. One conditioning factor and the
original normalized fan fraction are retained.

For every $0\le j<m$,
$\max(j,m-1-j)\ge(m-1)/2$.
Sum the nonnegative bounds and multiply by $m^2$ to obtain
\eqref{pc:all-marks}. Its scalar right-hand side tends to zero by
geometric decay and is independent of radius, labels and chosen
roof-value restrictions. Full positive sums precede every such
supremum. The local representatives were used only to prove the
finite-count inequalities, and no original outer limit is exchanged.
\end{proof}

\paragraph{What has been closed and what remains.}
Theorem~\ref{pc:fan} pays the complete \emph{leading} incoming
zero-clearance critical-fan coefficient in this actual physical
stratum, including positive corner fans with collapsed angular
slices and every shallow marked index. It extends the inherited
unscaled coefficient budget by an explicit collision-normalized
vanishing bound. It does not assert a uniform size for all local
critical coordinate neighborhoods, a collision-uniform sum of the
local coarea remainders, or a trace theorem at visited grazing.
An arbitrary Borel pairing selector need not have the analytic
sector property proved above; it may only restrict the already
bounded family when its own leading coefficient exists with weight
between zero and one.

The full regular angular marked charge is controlled by 320, and
the leading physical marked critical fan by this section.
Distinct visited-grazing, section/pairing boundary measures,
additional multipliers on those boundary measures, nonleading critical remainders
and the original same-roof positive-source comparison retain
their exact remaining obligations. In particular
\eqref{eq:sc-original-obligation} and the four-coordinate raw
inverse theorem are not asserted complete.
